% Drawing 73/01-10 -- folder 73, batch 01, archive page 1 (PDF page 2).
% Transcription: figure : $s_0$, $s_i$ et $s_{i+1} = s_{i+2}$ sur une même droite, notée $d_0 = \cdots = d_i$, puis la droite $d_{i+1}$ qui repart du point double vers le haut à droite.
% Kind: aligned vertices and the next side leaving the double point. Method 2 (schematic redraw from the facsimile). Status: ok.
% Notes: As on the page: the segment from s_0 through s_i to the heavy double point s_{i+1} = s_{i+2}, and the side d_{i+1} leaving that point; labels under the points, « d_0 = ... = d_i » under the segment, d_{i+1} beside the oblique side. Nothing left out. Redrawn 2026-11-12 (issue 128): to check again.
% Coordinates: 1 unit = 1 mm = 1/100 of the crop width (bbox in page fractions: [0.2, 0.53, 0.64, 0.635]).
% Redrawn by Claude (claude-fable-5-2), issue #128, 2026-11-12; not checked by a human; not compiled (no TeX in the session).
% Rights: reproduces part of the work -- behind RENDER_TRANSCRIPTS, never published by itself.
\documentclass[tikz,border=3pt]{standalone}
\usepackage{amsmath,amssymb}
\usetikzlibrary{patterns,arrows.meta}
\tikzset{pencil/.style={black!45}, struck/.style={black!35}, mpt/.style={circle,draw,inner sep=1.3pt}, pt/.style={circle,fill,inner sep=0.9pt}, opt/.style={circle,draw,fill=white,inner sep=1.1pt}}
\begin{document}
\begin{tikzpicture}[x=1mm,y=1mm,line cap=round,line join=round,font=\scriptsize]
\draw (18.2,-21) -- (55.1,-21.3) -- (73.8,-21.1) -- (93.9,-3.4);
\node[pt] at (18.2,-21) {};
\node[pt] at (55.1,-21.3) {};
\node[pt] at (73.8,-21.1) {};
\node[inner sep=1pt] at (16,-23.7) {\small $s_0$};
\node[inner sep=1pt] at (55.1,-24.4) {\small $s_i$};
\node[inner sep=1pt] at (80.5,-24.9) {\small $s_{i+1}=s_{i+2}$};
\node[inner sep=1pt] at (88.8,-15.3) {\small $d_{i+1}$};
\node[inner sep=1pt] at (37.9,-28) {\small $d_0=\cdots=d_i$};
\end{tikzpicture}
\end{document}
